% =====================================================================
%  CAPÍTULO 7 — TESTES E VALIDAÇÃO
% =====================================================================
\chapter{Testes e Validação}
\label{cap:testes}

A validação do sistema combina três níveis: testes automatizados do
\emph{backend}, verificação estática por tipagem, e validação funcional do
\emph{pipeline} completo. Discutem-se ainda as métricas relevantes para
uma avaliação académica da qualidade dos resultados.

\section{Plano de Testes}
\label{sec:plano_testes}

A estratégia de testes assentou em três pilares complementares:

\begin{enumerate}[leftmargin=*]
  \item \textbf{Testes unitários} — verificam, isoladamente, a lógica de
  componentes críticos e determinísticos.
  \item \textbf{Testes de integração} — verificam o comportamento de
  \emph{endpoints} da API, incluindo autenticação e sondas de saúde.
  \item \textbf{Verificação estática} — a tipagem (Python \emph{type
  hints} e TypeScript estrito) atua como uma camada de validação contínua,
  capturando incompatibilidades de esquema entre \emph{frontend} e
  \emph{backend} em tempo de compilação.
\end{enumerate}

\section{Testes Automatizados (pytest)}
\label{sec:pytest}

A suíte de testes do \emph{backend} foi implementada com \textbf{pytest}
(com \texttt{pytest-asyncio} e \texttt{pytest-cov}) e conta atualmente com
\textbf{37 testes a passar}. A \Cref{tab:testes} resume a sua distribuição.

\begin{table}[H]
\centering
\caption{Distribuição da suíte de testes automatizados.}
\label{tab:testes}
\small
\begin{tabularx}{\linewidth}{@{}llp{6.6cm}@{}}
\toprule
\textbf{Tipo} & \textbf{Módulo de teste} & \textbf{Âmbito} \\
\midrule
Unitário    & \texttt{test\_security}        & \emph{Security scanner}: \emph{checksum} e deteção de ficheiros inseguros. \\
Unitário    & \texttt{test\_upload\_validator} & Validação de tipo, \emph{magic bytes} e limites de tamanho. \\
Unitário    & \texttt{test\_date\_resolver}  & Resolução de datas incompletas e ambíguas (M2). \\
Unitário    & \texttt{test\_family\_graph}   & Construção do grafo e resumos de parentesco (M2). \\
Unitário    & \texttt{test\_templates}       & Integridade e formatação dos \emph{templates} narrativos (M3). \\
Unitário    & \texttt{test\_video\_builder}  & Funções de enquadramento e montagem de vídeo (M4). \\
Integração  & \texttt{test\_auth\_routes}    & Fluxo de autenticação na API. \\
Integração  & \texttt{test\_health\_route}   & Sondas de \emph{liveness} e \emph{deep health}. \\
\bottomrule
\end{tabularx}
\end{table}

A \Cref{fig:pytest} apresenta a execução da suíte. Os testes incidem
deliberadamente sobre a lógica \emph{determinística} (resolução de datas,
grafo, validação, \emph{templates}): a saída dos \textsc{llm} é
não-determinística e, por isso, validada por outras vias
(\Cref{sec:qualidade}).

\begin{figure}[H]
\centering
% \includegraphics[width=0.9\linewidth]{pytest.png}
\figph[4.5cm]{Captura de ecrã da execução \texttt{pytest -v} mostrando
os 37 testes a passar (idealmente com o sumário de cobertura
\texttt{--cov}).}
\caption{Resultado da execução da suíte de testes automatizados.}
\label{fig:pytest}
\end{figure}

\section{Validação Funcional do \emph{Pipeline}}
\label{sec:validacao_funcional}

Para além dos testes automatizados, validou-se o \emph{pipeline} completo
ponta-a-ponta com material de teste — incluindo ficheiros GEDCOM de
exemplo de várias famílias. O percurso \emph{happy path} validado foi:

\begin{enumerate}[leftmargin=*]
  \item Registo e autenticação de um utilizador.
  \item Carregamento de um conjunto de fotografias (M1: \textsc{exif} +
  análise visual + \emph{Storage}).
  \item Importação de uma árvore GEDCOM (M1/M2: pessoas + grafo).
  \item Geração de uma narrativa com tema, tom e idioma escolhidos (M3).
  \item Geração do vídeo documental a partir da narrativa (M4).
  \item Reprodução e descarregamento do vídeo final.
\end{enumerate}

Em cada etapa verificaram-se o estado das tarefas, as notificações ao
utilizador e a coerência dos artefactos produzidos.

\section{Métricas de Qualidade Narrativa}
\label{sec:qualidade}

A avaliação da qualidade da narrativa — intrinsecamente subjetiva — é um
desafio reconhecido (\Cref{sub:desafios}). Definem-se duas vias
complementares, em parte já instrumentáveis com os dados que o sistema
persiste (\texttt{task\_records}, \texttt{stories.facts\_used}):

\begin{itemize}[leftmargin=*]
  \item \textbf{Métricas objetivas.}
  \begin{itemize}
    \item \emph{Aderência ao PT-PT}: contagem automática de
    ``brasileirismos'' proibidos (a lista existe em \texttt{PT\_PT\_RULES})
    no texto gerado, produzindo uma percentagem de violações.
    \item \emph{Desempenho}: tempo médio de geração de narrativa e de
    vídeo, e taxa de sucesso/falha por \emph{template} e por \emph{backend}
    de \textsc{llm} (Ollama vs.\ Gemini).
  \end{itemize}
  \item \textbf{Métricas subjetivas.} Avaliação por um painel humano
  (escala 1--5) das dimensões \emph{coerência}, \emph{fidelidade factual}
  e \emph{qualidade literária}, sobre uma amostra de narrativas geradas.
\end{itemize}

A \Cref{tab:metricas} propõe a grelha de avaliação. A recolha sistemática
destes dados e a sua apresentação num painel constituem trabalho em curso,
discutido no \Cref{cap:conclusoes}.

\begin{table}[H]
\centering
\caption{Grelha proposta para avaliação da qualidade das narrativas
geradas.}
\label{tab:metricas}
\small
\begin{tabularx}{\linewidth}{@{}p{3.6cm}Yc@{}}
\toprule
\textbf{Métrica} & \textbf{Como medir} & \textbf{Tipo} \\
\midrule
Aderência PT-PT       & \% de violações às regras \texttt{PT\_PT\_RULES} & Objetiva \\
Fidelidade factual    & \% de factos da narrativa presentes no material   & Mista \\
Coerência             & Pontuação humana 1--5                              & Subjetiva \\
Qualidade literária   & Pontuação humana 1--5                              & Subjetiva \\
Tempo de geração      & Segundos (narrativa / vídeo)                       & Objetiva \\
Taxa de sucesso       & \% de gerações concluídas por \emph{template}      & Objetiva \\
\bottomrule
\end{tabularx}
\end{table}

\section{Resultados e Bugs Resolvidos}
\label{sec:resultados_testes}

Durante o desenvolvimento foram identificados e corrigidos diversos
defeitos relevantes, dos quais se destacam: a desalinhamento de nomes de
campos entre \emph{frontend} e \emph{backend} (por exemplo,
\texttt{narrative} vs.\ \texttt{content}, e \texttt{completed} vs.\
\texttt{ready} no estado dos vídeos), captado e resolvido graças à tipagem
estrita; a questão das tarefas eternamente ``em execução'', resolvida com
o \emph{sweep} de órfãs (\Cref{lst:sweep}); e a migração de
\texttt{datetime.utcnow()} para \texttt{datetime.now(UTC)} em conformidade
com as versões recentes do Python. Globalmente, o sistema cumpre os
requisitos funcionais essenciais (\rf{1}--\rf{13}), produzindo
narrativas e vídeos coerentes a partir de arquivos familiares reais.
